\section{Results}
\label{sec:results}

\subsection{Controlled Evidence-Contract Results}

Table~\ref{tab:mutation-results} reports every frozen candidate. All 20 valid
reports passed deterministic verification. All 16 mutated reports were
intercepted, yielding 20/20 valid-report acceptance, 16/16 hard-defect recall,
and 0/20 false blocks on this controlled benchmark. The verifier routed
repairable selector and grounding defects to \textsc{needs-evidence}; it routed
scope, hash-integrity, and diagnosis-conflict defects directly to
\textsc{review-required}.

\begin{table}[t]
    \caption{Deterministic verification on the 36-candidate frozen benchmark.
    ``NE'' denotes \textsc{needs-evidence}; ``HR'' denotes
    \textsc{review-required}.}
    \label{tab:mutation-results}
    \centering
    \small
    \begin{tabular}{@{}lrrr@{}}
        \toprule
        Candidate family & $n$ & NE & HR \\
        \midrule
        Valid, unmodified & 20 & 0 & 0 \\
        Unsupported scope & 6 & 0 & 6 \\
        Invalid selector & 2 & 2 & 0 \\
        Value/hash mismatch & 2 & 0 & 2 \\
        Claim not grounded & 2 & 2 & 0 \\
        Critical span not grounded & 2 & 2 & 0 \\
        Diagnosis conflict & 2 & 0 & 2 \\
        \midrule
        All injected defects & 16 & 6 & 10 \\
        \bottomrule
    \end{tabular}
\end{table}

These results test exact implemented rules against programmed defects. They do
not estimate performance on naturally occurring diagnoses and do not measure
the optional semantic verifier. Their contribution is narrower: the
claim--evidence contract creates an executable boundary that distinguishes
valid reports from six classes of known invalid report.

\subsection{System Validation Results}

Table~\ref{tab:system-results} summarizes the system-level evidence. The audited
public snapshot at Git revision
\nolinkurl{441871aa19cd4d7c129a721a449c5a098780afd1} collects 1,638 tests;
its complete local release gate recorded 1,637 passes, one platform-specific
skip, and 93.40\% statement coverage.
Static analysis, packaging, Compose validation, and the unprivileged container
path passed. The accepted recovery gate also completed all 20 independent
SQLite stability processes and preserved an end-to-end review decision across
an API-container restart.

\begin{table}[t]
    \caption{SpanVouch engineering validation. The exact public Git revision
    is stated in the accompanying text.}
    \label{tab:system-results}
    \centering
    \small
    \begin{tabular}{@{}p{0.30\columnwidth}p{0.62\columnwidth}@{}}
        \toprule
        Validation surface & Observed result \\
        \midrule
        Code snapshot & Public Git revision named above \\
        Public contracts & 6 versioned roots with strict valid/invalid fixtures \\
        Release suite & 1,637 passed, 1 skipped; 93.40\% coverage \\
        Static/delivery gates & Ruff, strict mypy, locked wheel, and Compose passed \\
        Process recovery & 20/20 SQLite processes completed \\
        Container recovery & Unprivileged API retained the terminal decision after restart \\
        Frozen artifact & Manifest and metrics hashes reproduced exactly \\
        \bottomrule
    \end{tabular}
\end{table}

\subsection{Offline Matrix Results}

The offline matrix evaluated all 24 scheduled domain--framework--condition
cells. It completed four adapter executions and both framework-parity
comparisons, then generated the evaluated-results manifest, statistical assets,
claim-gate report, and paper assets. Repeated executions produced identical
bundles. The run attempted zero provider requests and zero GPU operations.

The matrix result establishes that every B0--B5 delivery and analysis interface
is executable across both labs and framework adapters. It does not rank the six
conditions: fake semantic-provider outputs intentionally make the recorded
condition risks unsuitable for an effectiveness comparison. Real-provider
false-acceptance, correlated-error, and empirical target-risk results therefore
remain outside the reported findings.
